\begin{figure}[tb]
\centering
\begin{tikzpicture}
\begin{axis}[axis main,width=.82\linewidth,height=6.3cm,axis equal image,
 xlabel={$i_d$},ylabel={$i_q$},xmin=-3,xmax=2.2,ymin=-1.7,ymax=1.7]
 \addplot[voltage curve,BrickRed!45,domain=0:360,samples=121]
 ({-.25+1.38*cos(x)*cos(38)-.95*sin(x)*sin(38)},
  {1.38*cos(x)*sin(38)+.95*sin(x)*cos(38)});
 \addplot[voltage curve,BrickRed!72,domain=0:360,samples=121]
 ({-1.0+1.18*cos(x)*cos(16)-.63*sin(x)*sin(16)},
  {1.18*cos(x)*sin(16)+.63*sin(x)*cos(16)});
 \addplot[voltage curve,BrickRed,domain=0:360,samples=121]
 ({-1.55+.78*cos(x)},{.39*sin(x)});
 \node[annotation] at (axis cs:1.25,1.25){low speed:\\nearly circular};
 \node[annotation] at (axis cs:-.2,-1.25){medium};
 \node[annotation] at (axis cs:-2.15,.75){high speed:\\axes align};
\end{axis}
\end{tikzpicture}
\caption{Resistance and saliency generate the mixed term and therefore the
tilt. At low speed the formal axis angle is poorly meaningful because the
shape approaches a resistance circle; at high speed the axes approach $dq$.}
\label{fig:psm-tilt}
\end{figure}
